\documentclass[12pt,a4paper]{article}
\usepackage[utf8]{inputenc}
\usepackage[spanish]{babel}
\usepackage{amsmath, amssymb, amsfonts}
\usepackage{graphicx}
\usepackage{hyperref}
\usepackage{color}
\usepackage{geometry}
\usepackage{float}
\usepackage{booktabs}
\usepackage{multirow}
\usepackage{longtable}
\usepackage{fancyhdr}
\usepackage{titlesec}
\usepackage{enumitem}

% Configuración de márgenes
\geometry{
    left=2.5cm,
    right=2.5cm,
    top=2.5cm,
    bottom=2.5cm
}

% Configuración de encabezados
\pagestyle{fancy}
\fancyhf{}
\fancyhead[L]{\leftmark}
\fancyhead[R]{\thepage}
\renewcommand{\headrulewidth}{0.5pt}

% Configuración de títulos
\titleformat{\section}{\Large\bfseries}{\thesection}{1em}{}
\titleformat{\subsection}{\large\bfseries}{\thesubsection}{1em}{}
\titleformat{\subsubsection}{\normalsize\bfseries}{\thesubsubsection}{1em}{}

% Configuración de enlaces
\hypersetup{
    colorlinks=true,
    linkcolor=blue,
    urlcolor=blue,
    citecolor=blue
}

\begin{document}

\begin{titlepage}
    \begin{center}
        \vspace*{1cm}
        
        \Huge
        \textbf{PROYECTO DE INVESTIGACIÓN}
        
        \vspace{0.5cm}
        \LARGE
        Efectividad de los Juegos Gramaticales Bilingües en el\\
        Desarrollo de Competencias Morfosintácticas en\\
        Estudiantes de Español e Inglés
        
        \vspace{1.5cm}
        
        \textbf{Un Estudio Cuasi-Experimental}
        
        \vspace{2cm}
        
        \Large
        \textbf{Investigador Principal:} [Nombre del investigador]\\
        \textbf{Institución:} [Nombre de la universidad]\\
        \textbf{Departamento:} Lingüística Aplicada / Tecnología Educativa
        
        \vspace{2cm}
        
        \large
        \textbf{Duración:} 24 meses\\
        \textbf{Fecha de Inicio:} [Fecha]\\
        \textbf{Financiamiento Solicitado:} [Monto]
        
        \vspace{1cm}
        
        \textbf{Palabras Clave:} Gamificación, bilingüismo, adquisición de segundas lenguas, tecnología educativa, morfosintaxis, aprendizaje adaptativo
        
        \vfill
        
        \today
    \end{center}
\end{titlepage}

\newpage
\tableofcontents
\newpage

\section{Datos Generales del Proyecto}

\begin{table}[H]
\centering
\begin{tabular}{|l|p{10cm}|}
\hline
\textbf{Investigador Principal} & [Nombre del investigador] \\
\hline
\textbf{Institución} & [Nombre de la universidad/institución] \\
\hline
\textbf{Departamento} & Departamento de Lingüística Aplicada / Tecnología Educativa \\
\hline
\textbf{Duración} & 24 meses \\
\hline
\textbf{Fecha de Inicio} & [Fecha] \\
\hline
\textbf{Fecha de Finalización} & [Fecha] \\
\hline
\textbf{Financiamiento Solicitado} & [Monto en moneda local] \\
\hline
\textbf{Palabras Clave} & Gamificación, bilingüismo, adquisición de segundas lenguas, tecnología educativa, morfosintaxis, aprendizaje adaptativo \\
\hline
\end{tabular}
\end{table}

\section{Planteamiento del Problema}

\subsection{Contextualización}

En las últimas décadas, la globalización ha incrementado la demanda de competencias bilingües, especialmente en el par lingüístico español-inglés. Sin embargo, los métodos tradicionales de enseñanza de gramática enfrentan desafíos significativos:

\begin{itemize}
    \item \textbf{Baja motivación:} Los ejercicios gramaticales repetitivos generan desinterés en los estudiantes.
    \item \textbf{Desconexión contextual:} Las estructuras se enseñan aisladas de su uso real.
    \item \textbf{Falta de personalización:} Dificultad para adaptarse al ritmo individual de cada estudiante.
    \item \textbf{Transferencia limitada:} Los estudiantes no aplican lo aprendido en contextos comunicativos reales.
\end{itemize}

\subsection{Problema de Investigación}

¿En qué medida los juegos gramaticales bilingües basados en gamificación y aprendizaje adaptativo mejoran el desarrollo de competencias morfosintácticas en estudiantes de español e inglés, en comparación con métodos tradicionales de enseñanza?

\subsection{Preguntas Específicas}

\begin{enumerate}
    \item ¿Existen diferencias significativas en la adquisición de estructuras gramaticales entre estudiantes que utilizan la plataforma gamificada y aquellos que reciben instrucción tradicional?
    \item ¿Cómo influye el nivel de dificultad adaptativo en la motivación y el compromiso de los estudiantes?
    \item ¿Qué patrones de error son más frecuentes en cada nivel y cómo se relacionan con la lengua materna?
    \item ¿La gamificación (logros, rachas, puntuaciones) afecta la retención a largo plazo de conceptos gramaticales?
    \item ¿Existen diferencias en la efectividad según el grupo de edad (primaria, secundaria, adultos)?
\end{enumerate}

\section{Objetivos}

\subsection{Objetivo General}

Evaluar la efectividad de una plataforma de juegos gramaticales bilingües (español-inglés) basada en gamificación y aprendizaje adaptativo para el desarrollo de competencias morfosintácticas en estudiantes de diferentes niveles educativos.

\subsection{Objetivos Específicos}

\begin{enumerate}
    \item \textbf{Diseñar e implementar} 30 juegos gramaticales adaptativos con tres niveles de dificultad (primaria, secundaria, adultos).
    
    \item \textbf{Desarrollar} un sistema de seguimiento de progreso que registre métricas detalladas de aprendizaje (tiempo de respuesta, patrones de error, rachas, etc.).
    
    \item \textbf{Comparar} el rendimiento gramatical de grupos experimentales (plataforma gamificada) vs. grupos de control (métodos tradicionales).
    
    \item \textbf{Analizar} la relación entre elementos de gamificación (logros, puntuaciones, competencia) y la motivación intrínseca de los estudiantes.
    
    \item \textbf{Identificar} patrones de error específicos por nivel y lengua materna para optimizar el diseño instruccional.
    
    \item \textbf{Evaluar} la retención a largo plazo de conceptos gramaticales a los 3, 6 y 12 meses posteriores a la intervención.
    
    \item \textbf{Determinar} la usabilidad y aceptación de la plataforma por parte de estudiantes y docentes.
\end{enumerate}

\section{Marco Teórico}

\subsection{Teoría de la Adquisición de Segundas Lenguas (SLA)}

El proyecto se fundamenta en los siguientes modelos teóricos:

\begin{itemize}
    \item \textbf{Hipótesis del Input Comprensible} (Krashen, 1985): El aprendizaje ocurre cuando el estudiante está expuesto a input ligeramente superior a su nivel actual (i+1). Los juegos adaptativos ajustan dinámicamente la dificultad.
    
    \item \textbf{Modelo de Procesamiento del Input} (VanPatten, 1996): La atención guiada a formas gramaticales específicas facilita la adquisición. Los juegos dirigen la atención a estructuras concretas.
    
    \item \textbf{Hipótesis de la Interacción} (Long, 1996): La negociación de significado y la retroalimentación correctiva promueven la adquisición. La plataforma proporciona feedback inmediato.
\end{itemize}

\subsection{Gamificación en Educación}

\begin{itemize}
    \item \textbf{Taxonomía de Elementos de Juego} (Bedwell et al., 2012): Metas, reglas, feedback, competencia, desafío, fantasía.
    
    \item \textbf{Modelo MDA} (Hunicke et al., 2004): Mecánicas (puntos, niveles), Dinámicas (recompensas, progresión), Estéticas (experiencia del jugador).
    
    \item \textbf{Teoría de la Autodeterminación} (Ryan \& Deci, 2000): La gamificación satisface necesidades de autonomía, competencia y relación.
\end{itemize}

\subsection{Aprendizaje Adaptativo}

\begin{itemize}
    \item \textbf{Modelo de Conocimiento del Estudiante} (Corbett \& Anderson, 1995): Seguimiento de habilidades y ajuste de dificultad.
    
    \item \textbf{Teoría de la Carga Cognitiva} (Sweller, 1988): Los juegos adaptativos optimizan la carga cognitiva.
\end{itemize}

\subsection{Enfoque Contrastivo}

\begin{itemize}
    \item \textbf{Análisis Contrastivo} (Lado, 1957): Las diferencias entre L1 y L2 predicen dificultades de aprendizaje.
    \item \textbf{Teoría de la Interferencia Lingüística} (Weinreich, 1953): Identificación de falsos amigos y estructuras problemáticas.
\end{itemize}

\subsection{Estado del Arte}

\begin{table}[H]
\centering
\caption{Estudios relacionados con gamificación en aprendizaje de idiomas}
\begin{tabular}{|p{3cm}|p{4cm}|p{3cm}|p{3cm}|}
\hline
\textbf{Estudio} & \textbf{Hallazgos} & \textbf{Limitaciones} & \textbf{Aportación} \\
\hline
Reinhardt \& Sykes (2014) & Los juegos digitales mejoran vocabulario en L2 & Enfoque limitado a léxico & Base para juegos de vocabulario \\
\hline
Cornillie et al. (2012) & Feedback correctivo en juegos es efectivo & Muestra pequeña & Diseño de retroalimentación \\
\hline
Hung et al. (2018) & Gamificación aumenta motivación en EFL & Sin seguimiento longitudinal & Sistema de logros \\
\hline
\textbf{Este proyecto} & Enfoque integral (30 juegos, 3 niveles) & Por evaluar & Contribución original \\
\hline
\end{tabular}
\end{table}

\section{Hipótesis}

\subsection{Hipótesis Principal (H1)}

\textbf{Los estudiantes que utilizan la plataforma de juegos gramaticales bilingües mostrarán un incremento significativamente mayor en competencias morfosintácticas (español-inglés) en comparación con aquellos que reciben instrucción tradicional, medido mediante pruebas estandarizadas pre-post intervención.}

\subsection{Hipótesis Secundarias}

\begin{itemize}
    \item \textbf{H2:} La tasa de retención de conceptos gramaticales a los 6 meses será superior en el grupo experimental (≥30\% de diferencia).
    
    \item \textbf{H3:} Existe una correlación positiva entre el número de logros desbloqueados y la motivación autorreportada por los estudiantes.
    
    \item \textbf{H4:} Los patrones de error difieren significativamente según la lengua materna y el nivel de dificultad.
    
    \item \textbf{H5:} La satisfacción y usabilidad de la plataforma superará el umbral de aceptación (SUS > 70) en todos los grupos de edad.
\end{itemize}

\section{Metodología}

\subsection{Diseño de Investigación}

\textbf{Tipo de estudio:} Cuasi-experimental con grupos de control no equivalentes y medidas repetidas.

\textbf{Variables:}

\begin{table}[H]
\centering
\begin{tabular}{|l|l|l|}
\hline
\textbf{Variable} & \textbf{Tipo} & \textbf{Indicadores} \\
\hline
Uso de plataforma gamificada & Independiente & Tiempo de uso, juegos completados \\
\hline
Competencia gramatical & Dependiente & Puntuación en pruebas estandarizadas \\
\hline
Motivación & Dependiente & Escala de motivación intrínseca (IMI) \\
\hline
Retención & Dependiente & Pruebas diferidas (3, 6, 12 meses) \\
\hline
Edad/Nivel educativo & Moderadora & Primaria, secundaria, adultos \\
\hline
Lengua materna & Moderadora & Español, inglés, otras \\
\hline
\end{tabular}
\end{table}

\subsection{Participantes}

\subsubsection{Criterios de Inclusión}

\begin{itemize}
    \item Estudiantes de español/inglés como L2
    \item Edades: 6-12 (primaria), 13-17 (secundaria), 18+ (adultos)
    \item Acceso a dispositivo con conexión a internet
    \item Consentimiento informado firmado
\end{itemize}

\subsubsection{Tamaño Muestral}

Cálculo mediante G*Power:
\begin{itemize}
    \item Tamaño del efecto esperado: d = 0.5 (medio)
    \item Potencia: 0.80
    \item α = 0.05
    \item \textbf{n = 64 por grupo} (128 total por nivel)
\end{itemize}

\textbf{Muestra total}: 384 participantes (128 por nivel, divididos en 64 experimental + 64 control)

\begin{table}[H]
\centering
\caption{Distribución de la muestra}
\begin{tabular}{|l|c|c|c|}
\hline
\textbf{Nivel} & \textbf{Grupo Experimental} & \textbf{Grupo Control} & \textbf{Total} \\
\hline
Primaria & 64 & 64 & 128 \\
Secundaria & 64 & 64 & 128 \\
Adultos & 64 & 64 & 128 \\
\hline
\textbf{Total} & \textbf{192} & \textbf{192} & \textbf{384} \\
\hline
\end{tabular}
\end{table}

\subsection{Instrumentos}

\subsubsection{Pruebas de Competencia Gramatical}

\begin{table}[H]
\centering
\begin{tabular}{|l|l|l|}
\hline
\textbf{Prueba} & \textbf{Descripción} & \textbf{Administración} \\
\hline
Pre-test & 50 ítems (25 español, 25 inglés) & Semana 1 \\
Post-test 1 & 50 ítems (paralelo al pre-test) & Semana 12 \\
Post-test 2 & 50 ítems (retención) & Semana 24 \\
Post-test 3 & 50 ítems (retención largo plazo) & Semana 48 \\
\hline
\end{tabular}
\end{table}

\subsubsection{Escalas Psicrométricas}

\begin{itemize}
    \item \textbf{Intrinsic Motivation Inventory (IMI)} - Adaptado para contexto educativo
    \item \textbf{System Usability Scale (SUS)} - Para evaluar usabilidad
    \item \textbf{Technology Acceptance Model (TAM)} - Utilidad percibida y facilidad de uso
\end{itemize}

\subsubsection{Registros Automáticos de la Plataforma}

\begin{itemize}
    \item Tiempo por juego
    \item Número de intentos
    \item Patrones de error
    \item Rachas y logros
    \item Progresión entre niveles
\end{itemize}

\subsection{Procedimiento}

\begin{enumerate}
    \item \textbf{Fase 1: Preparación (Meses 1-3)}
    \begin{itemize}
        \item Adaptación de instrumentos
        \item Entrenamiento de investigadores
        \item Selección de centros educativos
        \item Consentimientos informados
    \end{itemize}
    
    \item \textbf{Fase 2: Pre-intervención (Mes 4)}
    \begin{itemize}
        \item Administración de pre-test
        \item Registro de datos demográficos
        \item Asignación a grupos
    \end{itemize}
    
    \item \textbf{Fase 3: Intervención (Meses 5-8)}
    \begin{itemize}
        \item \textbf{Grupo experimental:} 12 semanas de uso de la plataforma (2 sesiones/semana, 45 minutos)
        \item \textbf{Grupo control:} Instrucción tradicional con mismos contenidos gramaticales
    \end{itemize}
    
    \item \textbf{Fase 4: Post-intervención (Mes 9)}
    \begin{itemize}
        \item Administración de post-test 1
        \item Aplicación de escalas (IMI, SUS, TAM)
    \end{itemize}
    
    \item \textbf{Fase 5: Seguimiento (Meses 10-16)}
    \begin{itemize}
        \item Post-test 2 (mes 12)
        \item Post-test 3 (mes 16)
    \end{itemize}
    
    \item \textbf{Fase 6: Análisis y Difusión (Meses 17-24)}
    \begin{itemize}
        \item Análisis estadístico
        \item Redacción de informes
        \item Publicaciones y congresos
    \end{itemize}
\end{enumerate}

\subsection{Análisis de Datos}

\begin{table}[H]
\centering
\caption{Técnicas estadísticas a utilizar}
\begin{tabular}{|l|l|l|}
\hline
\textbf{Objetivo} & \textbf{Técnica} & \textbf{Software} \\
\hline
Comparación pre-post & ANOVA mixto 2x3 & SPSS/R \\
Efecto de variables moderadoras & ANCOVA & SPSS/R \\
Correlación motivación-logros & Correlación de Pearson & SPSS/R \\
Patrones de error & Análisis de clusters & Python/R \\
Usabilidad & Estadísticos descriptivos & SPSS \\
\hline
\end{tabular}
\end{table}

\subsubsection{Modelo Estadístico Principal}

\begin{equation}
Y_{ijk} = \mu + \alpha_i + \beta_j + (\alpha\beta)_{ij} + \gamma_k + \epsilon_{ijk}
\end{equation}

Donde:
\begin{itemize}
    \item $Y_{ijk}$: Puntuación gramatical
    \item $\alpha_i$: Grupo (experimental/control)
    \item $\beta_j$: Tiempo (pre/post1/post2/post3)
    \item $(\alpha\beta)_{ij}$: Interacción grupo × tiempo
    \item $\gamma_k$: Nivel educativo (covariable)
    \item $\epsilon_{ijk}$: Error aleatorio
\end{itemize}

\section{Cronograma}

\begin{table}[H]
\centering
\caption{Cronograma de actividades por mes}
\begin{tabular}{|l|c|c|c|c|c|c|c|c|}
\hline
\textbf{Actividad} & \textbf{M1-3} & \textbf{M4} & \textbf{M5-8} & \textbf{M9} & \textbf{M10-12} & \textbf{M13-16} & \textbf{M17-20} & \textbf{M21-24} \\
\hline
Preparación & ███ & & & & & & & \\
Pre-test & & █ & & & & & & \\
Intervención & & & ████ & & & & & \\
Post-test 1 & & & & █ & & & & \\
Post-test 2 & & & & & ███ & & & \\
Post-test 3 & & & & & & ████ & & \\
Análisis & & & & & & & ████ & \\
Difusión & & & & & & & & ████ \\
\hline
\end{tabular}
\end{table}

\section{Presupuesto}

\begin{table}[H]
\centering
\caption{Presupuesto detallado del proyecto}
\begin{tabular}{|l|l|c|}
\hline
\textbf{Concepto} & \textbf{Descripción} & \textbf{Costo (USD)} \\
\hline
\multicolumn{3}{|l|}{\textbf{Personal}} \\
\hline
Investigador principal & 24 meses, dedicación parcial & 24,000 \\
Asistente de investigación & 12 meses & 12,000 \\
Programador & Desarrollo y mantenimiento & 8,000 \\
\hline
\multicolumn{3}{|l|}{\textbf{Equipamiento}} \\
\hline
Servidores cloud & 24 meses & 2,400 \\
Licencias software & SPSS, G*Power & 800 \\
Tabletas para evaluación & 20 unidades & 6,000 \\
\hline
\multicolumn{3}{|l|}{\textbf{Materiales}} \\
\hline
Impresión de instrumentos & & 500 \\
Incentivos participantes & & 2,000 \\
\hline
\multicolumn{3}{|l|}{\textbf{Difusión}} \\
\hline
Inscripción congresos & 2 congresos internacionales & 1,200 \\
Publicación open access & 2 artículos & 3,000 \\
\hline
\textbf{Imprevistos} & 10\% & 6,000 \\
\hline
\textbf{TOTAL} & & \textbf{65,900} \\
\hline
\end{tabular}
\end{table}

\section{Resultados Esperados}

\subsection{Productos Esperados}

\begin{table}[H]
\centering
\begin{tabular}{|l|l|l|}
\hline
\textbf{Producto} & \textbf{Descripción} & \textbf{Fecha estimada} \\
\hline
Plataforma funcional & 30 juegos implementados y validados & Mes 8 \\
Artículo 1 & Diseño y desarrollo de juegos gramaticales & Mes 12 \\
Artículo 2 & Efectividad de la gamificación en gramática & Mes 18 \\
Artículo 3 & Patrones de error y aprendizaje adaptativo & Mes 22 \\
Tesis doctoral & Del investigador principal & Mes 24 \\
Manual para educadores & Guía de implementación en aula & Mes 20 \\
Repositorio open source & Código de la plataforma & Mes 6 \\
\hline
\end{tabular}
\end{table}

\subsection{Impactos Esperados}

\subsubsection{Impacto Científico}
\begin{itemize}
    \item Contribución a la teoría de adquisición de segundas lenguas mediada por tecnología
    \item Validación empírica de gamificación en contextos bilingües
    \item Identificación de patrones de error por nivel y lengua materna
\end{itemize}

\subsubsection{Impacto Educativo}
\begin{itemize}
    \item Herramienta gratuita para instituciones educativas
    \item Metodología replicable para otros pares lingüísticos
    \item Reducción de brechas en aprendizaje de idiomas
\end{itemize}

\subsubsection{Impacto Social}
\begin{itemize}
    \item Democratización del acceso a herramientas de aprendizaje
    \item Fomento del bilingüismo en comunidades desfavorecidas
    \item Desarrollo de competencias para la empleabilidad
\end{itemize}

\section{Limitaciones y Riesgos}

\begin{table}[H]
\centering
\begin{tabular}{|l|l|l|l|}
\hline
\textbf{Riesgo} & \textbf{Probabilidad} & \textbf{Impacto} & \textbf{Mitigación} \\
\hline
Deserción de participantes & Media & Alto & Sobremuestreo inicial (20\%) \\
Problemas técnicos plataforma & Baja & Alto & Pruebas exhaustivas, soporte técnico \\
Contaminación entre grupos & Media & Medio & Centros educativos separados \\
Efecto novedad & Alta & Medio & Seguimiento longitudinal \\
Diferencias docentes & Media & Medio & Estandarización de instrucción \\
Acceso a tecnología & Baja & Alto & Provisión de dispositivos \\
\hline
\end{tabular}
\end{table}

\section{Aspectos Éticos}

\subsection{Consentimiento Informado}
\begin{itemize}
    \item Información clara sobre objetivos y procedimientos
    \item Confidencialidad de datos
    \item Derecho a retirarse en cualquier momento
\end{itemize}

\subsection{Protección de Datos}
\begin{itemize}
    \item Cumplimiento de GDPR / Ley de Protección de Datos
    \item Datos anonimizados para análisis
    \item Almacenamiento seguro en servidores institucionales
\end{itemize}

\subsection{Beneficencia y No Maleficencia}
\begin{itemize}
    \item La intervención no interfiere con currículo regular
    \item El grupo control recibe instrucción tradicional equivalente
    \item Todos los participantes acceden a la plataforma al finalizar
\end{itemize}

\subsection{Aprobación Ética}
\begin{itemize}
    \item Solicitud al Comité de Ética de la Investigación
    \item Certificado de buenas prácticas
    \item Permisos de instituciones educativas
\end{itemize}

\section{Referencias Bibliográficas}

\begin{thebibliography}{99}

\bibitem{bedwell2012}
Bedwell, W. L., et al. (2012). Toward a taxonomy linking game attributes to learning. \textit{Simulation \& Gaming}, 43(6), 729-760.

\bibitem{corbett1995}
Corbett, A. T., \& Anderson, J. R. (1995). Knowledge tracing: Modeling the acquisition of procedural knowledge. \textit{User Modeling and User-Adapted Interaction}, 4(4), 253-278.

\bibitem{cornillie2012}
Cornillie, F., et al. (2012). The role of feedback in foreign language learning through digital role playing games. \textit{Procedia-Social and Behavioral Sciences}, 34, 49-53.

\bibitem{deci2000}
Deci, E. L., \& Ryan, R. M. (2000). The "what" and "why" of goal pursuits: Human needs and the self-determination of behavior. \textit{Psychological Inquiry}, 11(4), 227-268.

\bibitem{deterding2011}
Deterding, S., et al. (2011). From game design elements to gamefulness: Defining "gamification". \textit{Proceedings of the 15th International Academic MindTrek Conference}, 9-15.

\bibitem{ellis2008}
Ellis, R. (2008). \textit{The Study of Second Language Acquisition}. Oxford University Press.

\bibitem{hung2018}
Hung, H. T., et al. (2018). Gamification in the foreign language classroom: A systematic review. \textit{Educational Technology \& Society}, 21(3), 104-118.

\bibitem{hunicke2004}
Hunicke, R., et al. (2004). MDA: A formal approach to game design and game research. \textit{Proceedings of the AAAI Workshop on Challenges in Game AI}.

\bibitem{krashen1985}
Krashen, S. D. (1985). \textit{The Input Hypothesis: Issues and Implications}. Longman.

\bibitem{lado1957}
Lado, R. (1957). \textit{Linguistics Across Cultures}. University of Michigan Press.

\bibitem{long1996}
Long, M. H. (1996). The role of the linguistic environment in second language acquisition. \textit{Handbook of Second Language Acquisition}, 413-468.

\bibitem{pane2014}
Pane, J. F., et al. (2014). \textit{Continued Progress: Promising Evidence on Personalized Learning}. RAND Corporation.

\bibitem{reinhardt2014}
Reinhardt, J., \& Sykes, J. M. (2014). Digital game and play activity in L2 teaching and learning. \textit{Language Learning \& Technology}, 18(2), 2-8.

\bibitem{sweller1988}
Sweller, J. (1988). Cognitive load during problem solving: Effects on learning. \textit{Cognitive Science}, 12(2), 257-285.

\bibitem{vanpatten1996}
VanPatten, B. (1996). \textit{Input Processing and Grammar Instruction}. Ablex Publishing.

\bibitem{vygotsky1978}
Vygotsky, L. S. (1978). \textit{Mind in Society}. Harvard University Press.

\bibitem{weinreich1953}
Weinreich, U. (1953). \textit{Languages in Contact}. Linguistic Circle of New York.

\end{thebibliography}

\section{Anexos}

\subsection{Anexo A: Prueba de Competencia Gramatical (Ejemplo)}

\textbf{Nivel Primaria - Español}
\begin{enumerate}
    \item El niño \_\_\_ (comer) una manzana.
    \item \_\_\_ perro es marrón. (El/La/Los)
    \item María y yo \_\_\_ (ser) amigos.
\end{enumerate}

\textbf{Nivel Primaria - Inglés}
\begin{enumerate}
    \item The cat \_\_\_ (sleep) on the sofa.
    \item \_\_\_ are my books. (This/These)
    \item She \_\_\_ (have) a red car.
\end{enumerate}

\subsection{Anexo B: Escala de Motivación Intrínseca (IMI) - Adaptada}

\begin{enumerate}
    \item Disfruté mucho usando esta plataforma.
    \item Los juegos me parecieron interesantes.
    \item Me sentí competente mientras jugaba.
    \item Me esforcé por hacerlo bien.
    \item Recomendaría esta plataforma a otros estudiantes.
\end{enumerate}

\subsection{Anexo C: Carta de Consentimiento Informado}

[Incluir modelo de carta con todos los elementos éticos requeridos]

\section{Firmas}

\begin{table}[H]
\centering
\begin{tabular}{|l|l|l|l|}
\hline
\textbf{Rol} & \textbf{Nombre} & \textbf{Firma} & \textbf{Fecha} \\
\hline
Investigador Principal & & & \\
Director de Tesis & & & \\
Comité de Ética & & & \\
Institución Educativa & & & \\
\hline
\end{tabular}
\end{table}

\section{Resumen Ejecutivo}

\begin{quote}
\textbf{Título:} Efectividad de los Juegos Gramaticales Bilingües en el Desarrollo de Competencias Morfosintácticas

\textbf{Objetivo:} Evaluar una plataforma de 30 juegos adaptativos para el aprendizaje de gramática español-inglés.

\textbf{Metodología:} Estudio cuasi-experimental con 384 participantes (primaria, secundaria, adultos), grupos experimental (gamificación) y control (tradicional), con medidas pre-post y seguimiento a 12 meses.

\textbf{Resultados esperados:} Mejora significativa en competencias gramaticales, alta motivación y retención, identificación de patrones de error.

\textbf{Impacto:} Contribución científica, herramienta educativa gratuita, democratización del aprendizaje de idiomas.

\textbf{Presupuesto:} USD 65,900 para 24 meses.
\end{quote}

\textit{Este proyecto de investigación ha sido diseñado siguiendo los estándares del Ministerio de Ciencia e Innovación y las directrices éticas internacionales para investigación con seres humanos.}

\end{document}
